% Full proofs for the GPSS complete-order boundary authoring unit.

\section{Future-only limits and complete-order boundary events}
\label{sec:annex-p09-gpss-complete-order-boundaries}

\begin{annexitem}{annex:p9.gpss.complete.order.bridge}{thm:p09-gpss-complete-order-bridge}{Order completeness and the future syntax bridge}
\begin{proof}
Work with a nonempty strict linear order. First suppose that every
nonempty subset bounded both above and below has an infimum and
supremum. Let \(A\) be nonempty and bounded below by \(l\).
Choose \(a_0\in A\) and put
\[
                         A_0=\{a\in A:a\leq a_0\}.
\]
This subset is nonempty, bounded below by \(l\), and bounded above
by \(a_0\). Let \(c=\inf A_0\). Since \(a_0\in A_0\),
\(c\leq a_0\). Every member of \(A\setminus A_0\) is above
\(a_0\) and hence above \(c\). Thus \(c\) is a lower bound
for all of \(A\). Any lower bound for \(A\) is a lower bound
for \(A_0\), so it is at most \(c\). This proves that \(c\)
is the infimum of \(A\).

For a nonempty upper-bounded \(A\), choose \(a_0\in A\) and
apply the two-sided bounded hypothesis to
\(A_1=\{a\in A:a_0\leq a\}\). Its supremum is an upper
bound for \(A\), since every discarded member is below \(a_0\).
Every upper bound for \(A\) bounds \(A_1\), proving leastness.
The converse follows by applying the two one-sided conditions to
a nonempty set with both bounds. This proves the equivalence with
the displayed \(\mathsf{Com}\) convention in the Kamp presentation.
The proof neither adds an endpoint nor assumes one exists.

For the temporal bridge, map corresponding unary letters to each other,
commute with Boolean operations, and map Kamp's \(U(C,B)\) to
\(B\mathbin{\mathsf U_s}C\). Their satisfaction clauses both say:
there is a point strictly after the current point where \(C\) holds,
and \(B\) holds at every intervening point. The witness is excluded
from the interval requirement, as is the current point. Induction
on the input formula therefore proves the bridge at every point
and for every unary interpretation. Quantifier renaming in the
first-order images preserves the same free current-point variable.

In the nested word formula, the induction starts with its last point
label. The step conjoins the current point label and uses one strict
future witness for the remainder, with the indicated open gap.
Expanding successively gives exactly the increasing finite realization
of the word. If an occurrence label is itself a supplied future
formula, the same induction substitutes its whole-model truth at each
chosen point; its internal witnesses need not remain in the word's
interval. Capture-free first-order copying yields the equivalent
label translation. No new unary interpretation is required.

This is the future part of the basic-word construction, with the two
sources' argument orders related explicitly. The proof has no rule
for mapping a since occurrence to \(L(U)\). In particular the
more general elementary class used by Kamp's later complete-order
induction is not identified with the permitted GPSS future labels.
\end{proof}
\end{annexitem}

\begin{annexitem}{annex:p9.gpss.right.limit}{thm:p09-gpss-right-limit}{Nonempty right accumulation with all endpoint cases}
\begin{proof}
For any current point \(x\), strict until gives
\[
 N^+A(x)\ \longleftrightarrow\
        \exists b\bigl(x<b\land
                    \forall z(x<z<b\to A(z))\bigr).
\]
Consequently classical negation of this formula with \(A\) replaced
by \(\neg A\) gives
\[
 \partial^+A(x)\ \longleftrightarrow\
        \forall b\bigl(x<b\to
                    \exists z(x<z<b\land A(z))\bigr).
\]
The formula \(F\top\) says exactly that a future point exists.
Their conjunction is the claimed nonempty limit condition.
These expansions prove that the condition concerns the predicate
at the quantified witness \(z\); it is independent of a change
at the current point alone when \(A\) is an atomic predicate.
For a compound \(A\), the assertion still makes no direct
current-point test of \(A(x)\).

If \(b\) is an immediate successor of \(x\), the interval
\((x,b)\) is empty. It satisfies the universal \(\neg A\)
condition, so both limit formulas are false. If \(x\) is last,
there is no future cut. The universal formula for \(\partial^+A\)
is true and \(F\top\) is false. This proves the stated endpoint
values, including a singleton order.

Now fix an actual \(y>x\) and write
\(R=\{z:x<z<y,\ A(z)\}\). If \(K^+A(x)\) holds,
the cut \(y\) supplies a member of \(R\). For every \(b>x\),
apply the limit condition to the smaller of \(b,y\); its witness
belongs to \(R\) and lies below \(b\). Hence \(R\) is
coinitial at \(x\).

Conversely, suppose \(R\ne\varnothing\) and \(x=\inf R\).
Every \(b>x\) fails to be a lower bound of \(R\), for otherwise
the greatest lower bound would be at least \(b>x\). There is
therefore \(z\in R\) with \(z<b\). Since \(z>x\), this is the
required strict witness. The point \(y>x\) supplies \(F\top\).
Equivalently one may define coinitiality by this very strict-witness
condition; the preceding argument proves its infimum formulation
without an assumption of density. All occurrences of \(A(z)\)
retain their original whole-model meaning.
\end{proof}
\end{annexitem}

\begin{annexitem}{annex:p9.gpss.future.word.limit}{thm:p09-gpss-future-word-limit}{A future-only finite-chain splice}
\begin{proof}
First, the starting formula \(E^+_W\) expresses a realization
beginning at the current point. For a one-point word this is the
point label. For a longer word, the strict until witness is its
second point, the gap from the current point has the first gap
label, and the suffix formula realizes all remaining positions.
Conversely every realization supplies exactly that witness and
suffix. Induction on the number of gaps proves the assertion,
including empty intervening intervals.

Let \(C_W(x)\) denote the body's nonempty right-accumulation
condition for a word \(W\). We prove
\[
 C_{(A)}(x)\longleftrightarrow K^+A(x)
\]
and, when \(W\) has \(n\geq1\) gaps,
\[
 C_W(x)\longleftrightarrow
                      K^+(E^+_W)(x)\land C_{W^-}(x).
 \tag{1}
\]
The base is the preceding right-limit lemma. For the forward
direction of (1), fix any \(b>x\). A realization wholly in
\((x,b)\) has its first point there and satisfies \(E^+_W\)
at that point. Deleting its first point leaves a realization of
\(W^-\) wholly in the same interval. The assumed future
nonemptiness remains true. Applying this to every \(b\) proves
both right-hand conditions.

For the converse, assume those two conditions and fix \(b>x\).
Choose a suffix realization
\[
                     x<v_1<\cdots<v_n<b.
\]
Its position \(v_i\) satisfies \(A_i\), and its gaps satisfy
\(B_i\) for \(2\leq i\leq n\). When \(n=1\) it is
just the point \(v_1\). Apply \(K^+(E^+_W)(x)\) to the
actual cut \(v_1>x\). It gives \(u_0\) with
\(x<u_0<v_1\) satisfying \(E^+_W\). Choose a whole-word
realization
\[
                         u_0<u_1<\cdots<u_n
\]
beginning there. Its later points are not initially required to
lie below \(b\).

If \(u_i<v_i\) for every \(1\leq i\leq n\), then
\(u_n<v_n<b\); this whole realization already lies in
\((x,b)\). Otherwise let \(k\) be the least index
\(1\leq k\leq n\) with \(v_k\leq u_k\).
If \(k=1\), the choice of \(u_0\) gives \(u_0<v_1\).
If \(k>1\), minimality gives
\(u_{k-1}<v_{k-1}<v_k\). In either case
\[
                         u_{k-1}<v_k\leq u_k.
\]
Form the chain
\[
                   u_0,\ldots,u_{k-1},v_k,\ldots,v_n.
 \tag{2}
\]
It has exactly \(n+1\) strictly increasing points. It starts
above \(x\), ends below \(b\), and each position has the
required \(A_i\) label. The gaps before index \(k\) are old
whole-word gaps; those after \(k\) are suffix gaps. For the
only new gap,
\[
              (u_{k-1},v_k)\subseteq(u_{k-1},u_k),
\]
so the old gap supplies \(B_k\). Equality \(v_k=u_k\)
causes no difficulty and introduces no point test for that gap.
This verifies every condition on (2).

The case \(k=1\) replaces the entire old suffix. The case
\(k=n\) replaces only its last point. When \(n=1\) the
same argument joins the chosen \(u_0\) to \(v_1\) unless
the old endpoint is already earlier. Thus no exceptional word
length or extreme crossing is omitted. Since \(b>x\) was
arbitrary, the right-hand side of (1) yields \(C_W(x)\).
Its future nonemptiness follows from either \(K^+\) factor.

Induction on \(n\), using (1), now proves that \(\Gamma_W\)
expresses \(C_W\). The recursion in the body constructs both
the starting formulas and the limit formulas by finite syntax
operations. All labels remain in \(L(U)\), since their only
added modalities are future until and its displayed Boolean
abbreviations.

For the exact bridge to Kamp's limiting condition, remove only
the existence of a future point. At a last point the universal
condition over all right cuts is true. At any other point it is
exactly \(C_W\). Thus its formula is
\(\neg F\top\lor\Gamma_W\). For \(n\geq1\), choosing
the first and last positions of the realization gives Kamp's
strictly ordered \(u,v\) inside every right interval, and
conversely that basic-pattern realization gives all the positions.
At \(n=0\), the single point gives \(\partial^+A\).
This proves the exact semantic bridge to the earlier
\(\operatorname{Lf}_n\) grammar.

The displayed earlier future-limit recurrence uses a formula for
a realization ending at a point; that formula can contain since.
Equation (1) instead uses only a formula for a realization
starting at a point and the proved splice (2). Its future-only
conclusion follows from this new derivation, not from suppressing
a since symbol in the earlier construction.

For the example \(W=(A,B,C)\), substituting the definitions
gives precisely the two displayed \(K^+\) factors. If
\(A=(0,1)\), \(B=\mathbb R\), and \(C=\{2\}\),
every right neighborhood of zero has an \(A\)-point with a
future \(C\)-witness at 2. The first factor holds. No sufficiently
small neighborhood has any \(C\)-point, so the second factor
and the required word limit fail. With \(A=C=(0,1)\),
every neighborhood of zero contains two such increasing points
and their open gap satisfies \(B\). Both the word condition
and the formula are true.
\end{proof}
\end{annexitem}

\begin{annexitem}{annex:p9.gpss.fixed.start.limit}{thm:p09-gpss-fixed-start-limit}{The retained initial point and its open gap}
\begin{proof}
Suppose the nonempty fixed-start accumulation condition holds.
A future point exists, so at least one of its required
realizations exists. Its starting label gives \(A_0(x)\).
Its first internal point \(z_1>x\) supplies a strict
right neighborhood \((x,z_1)\) on which \(B_1\) holds.
Hence \(N^+B_1(x)\). For every right cut, deleting the
fixed initial point from a realization leaves the suffix word
wholly in that interval. The future-word-limit lemma gives
\(\Gamma_{W^-}(x)\).

Conversely assume the three displayed factors. Choose \(d>x\)
with \(B_1\) throughout \((x,d)\), using \(N^+B_1\).
For an arbitrary \(b>x\), use \(\Gamma_{W^-}\) to realize
that suffix wholly in \((x,\min(b,d))\). This contains at
least its first point \(z_1\). Prepend \(x\).
The label \(A_0\) holds there, and
\((x,z_1)\subseteq(x,d)\) supplies \(B_1\).
Every subsequent point and gap comes from the suffix realization.
The resulting whole word starts at \(x\) and ends below \(b\).
Existence of \(d\) also gives a future point. This proves both
directions, including the one-gap case.

If the existence requirement is omitted, a last point makes the
universal fixed-start condition vacuous. The separate disjunct
\(\neg F\top\) handles it, regardless of \(A_0(x)\).
At a nonlast point the preceding equivalence applies. If an
immediate successor exists, the interval to that successor
contains no suffix point; \(\Gamma_{W^-}\) is false, as is
the fixed-start accumulation condition.
\end{proof}
\end{annexitem}

\begin{annexitem}{annex:p9.gpss.left.labelled.closure}{thm:p09-gpss-left-labelled-closure}{Positive assembly and exact future-ray complements}
\begin{proof}
For two left-labelled decompositions at the same endpoints,
write them as \(\alpha(x)\land D(x,y)\) and
\(\beta(x)\land E(x,y)\). Their conjunction is
\[
           (\alpha\land\beta)(x)\land(D(x,y)\land E(x,y)).
\]
The proved positive decomposition construction gives a finite
list for the second factor, including coincident interior
witnesses and every transferred gap condition. Prefix each
output by the same future label \(\alpha\land\beta\).
Distribute finite input disjunctions by their finite Cartesian
product. The empty product gives left label true and the
zero-point true-gap decomposition; a false factor gives an
empty list. This proves all finite conjunction cases.

For a guarded concatenation at \(z\), a left constituent is
\(\alpha(a)\land D(a,z)\), and a right constituent is
\(\beta(z)\land E(z,b)\). If \(\theta(z)\) is the
required midpoint label, conjoin \(\beta\land\theta\)
at that point in the ordinary concatenation. Retain
\(\alpha(a)\) as the final left label. This gives the
original existential with the same midpoint witness in both
directions. Distributing over the finite input lists gives
the general construction. The lower endpoint \(a\) remains
actual, and each quantified \(z\) is an actual interior
point. No label is placed at a formal endpoint.

On a ray, the suffix-translation lemma converts each original
decomposition \(D(x,+\infty)\) to its nested strict-until
formula, including its terminal \(G\) condition. Each global
cofinal conjunct has its already proved exact formula
\(F\Psi\land GF\Psi\). Conjoin the supplied left label
and disjoin the finite list to obtain \(B\in L(U)\).
An empty list translates to false.

The decomposition \(D_{\top}(x,+\infty)\) has no witness
points and requires only truth at all later points, so it is
true at every \(x\). In particular it is true at a last point
by vacuity. Therefore
\((\neg B)(x)\land D_{\top}(x,+\infty)\) is exactly
the complement of the given list at every point. Its sole
nontrivial condition is an allowed lower-endpoint label.
The construction does not need a successor and works on
arbitrary linear orders. It does not return to the narrower
endpoint-free class.
\end{proof}
\end{annexitem}

\begin{annexitem}{annex:p9.gpss.infimum.descriptors}{thm:p09-gpss-infimum-descriptors}{Empty, initial, attained and unattained infima}
\begin{proof}
Fix \(s<t\) and the set \(R=A_{s,t}\). It is empty exactly
when the first displayed universal condition holds. Suppose
it is nonempty. Completeness supplies \(c=\inf R\), since
\(s\) is a lower bound. Every \(r\in R\) satisfies
\(c\leq r<t\), whence \(c<t\), and greatestness of the
lower bound gives \(s\leq c\).

If \(c=s\), the right-limit lemma gives \(K^+A(s)\),
because \(R\) is a nonempty subset of \((s,t)\) with
infimum \(s\). Conversely that formula gives precisely
this infimum and nonemptiness, by the same lemma. Since
\(s\notin(s,t)\), its own membership in the interpretation
of \(A\) plays no part. The initial status is therefore
neither a test of \(A(s)\) nor an attained minimum of \(R\).

Suppose \(s<c<t\). If \(A(c)\), then \(c\in R\), so
it is the minimum. No point of \(R\) lies below \(c\);
this is the absence condition on \((s,c)\). Conversely,
\(A(c)\) and that absence, with the strict bounds, put
\(c\) in \(R\) and below or equal to every member, so
it is the attained infimum.

If \(A(c)\) is false, \(R\) has no member at or below
\(c\). For any \(d>c\), greatestness of the lower bound
implies that some \(r\in R\) satisfies \(r<d\);
otherwise \(d\) would be a larger lower bound. This
member satisfies \(c<r<d\). A future point exists,
for example \(t>c\), so \(K^+A(c)\) follows.
The absence on \((s,c)\) and the false point test at
\(c\) are thus necessary.

Conversely, suppose the unattained row holds. Every point
of \(R\) is strictly above \(c\), by the absence
condition and \(\neg A(c)\). Apply \(K^+A(c)\) to
the cut \(t\) to obtain a member of \(R\), so the set
is nonempty. For any \(d>c\), apply it to \(\min(d,t)\)
to get a member of \(R\) below \(d\). No such \(d\)
can be a lower bound. Thus \(c\) is the infimum and
does not belong to \(R\).

Infima are unique in a strict linear order: if \(c,d\)
are both greatest lower bounds, each is at most the other.
The three defined statuses are mutually exclusive and,
together with emptiness, exhaust the possibilities just proved.
Only their exhaustion from nonemptiness required completeness;
each described value is verified as an infimum directly.
If \((s,t)\) is empty, its occurrence set is empty,
regardless of the label, and none of the defined cases can hold.

For the future ray put \(R=\{u:s<u,\ A(u)\}\).
Its emptiness is \(G\neg A(s)\). If it is nonempty,
it is bounded below by \(s\) and hence has an infimum
\(c\geq s\), by the complete-order bridge. If \(c=s\),
the same greatest-lower-bound argument gives precisely
\(K^+A(s)\). If \(c>s\) is attained, absence of
earlier occurrences and the point label \(A(c)\)
give \((\neg A)\mathbin{\mathsf U_s}A\) at \(s\).
Conversely its witness is a first member of the
future occurrence set. If \(c>s\) is unattained,
absence before \(c\), falsehood at \(c\), and
arbitrary occurrences immediately after \(c\)
give the fourth displayed formula. Conversely
its witness and \(K^+A\) make it exactly this
unattained infimum, by the same lower-bound
argument as in the bounded case. The three
nonempty cases have a unique infimum and distinct
membership or endpoint status; they cannot
overlap. At a last point the set is empty,
\(G\neg A\) is true, and all three other
formulas are false.

The four real-order examples follow by taking their ordinary
infima: the empty set has none, \((0,1)\) has infimum
0 outside the interval, \(\{1\}\) has its attained
minimum at 1, and \((1,3/2)\) has an unattained infimum
at 1. Changes at the excluded initial point do not affect
these sets when \(A\) is the displayed atomic interpretation.
\end{proof}
\end{annexitem}

\begin{annexitem}{annex:p9.gpss.infimum.case.compiler}{thm:p09-gpss-infimum-case-compiler}{Every finite infimum-event diagram returns to an explicit formula}
\begin{proof}
The finite list has a topological ordering: an event's arguments
are roots or earlier values. Induction in that ordering shows
that every defined nonroot value lies in \([x,y)\).
For an infimum, its valid lower bound is at least \(x\),
and any occurrence witnessing nonemptiness is strictly below
its upper bound, which is at most \(y\). Its infimum is
therefore at least \(x\) and strictly below \(y\).
A test returns its argument, which is \(x\) or an earlier
nonroot value and hence also below \(y\). This explains
both the allowance of equality with \(x\) and the exclusion
of a test at the root \(y\).

Enumerate a defined/undefined choice for each event, one of
the three statuses for every defined infimum, and an ordered
partition of the roots and all defined events into equality
classes. The roots are defined, \(x<y\), every other
defined class lies between or at \(x\) and strictly below
\(y\), and all equality and strict comparisons are thereby
decided. There are finitely many such assignments. Reject
a case if its proposed statuses or order disagree with a
formal constraint described below. Keep only cases satisfying
the Boolean condition specified by the input.

For an infimum event with an undefined argument, require
the event to be undefined and add no occurrence constraint.
The same rule applies to defined but non-strictly ordered
bounds. These are consequences of the event's definition.
For defined ordered bounds \(s<t\), use exactly one
of the four descriptor rows proved in the preceding item.
An undefined output adds absence of \(A\) on \((s,t)\).
An initial output must equal \(s\) and adds \(K^+A\)
there. An attained output must lie strictly between the
bounds; it adds \(A\) at its class and absence of \(A\)
on \((s,c)\). An unattained output has the same strict
bounds, adds \(\neg A\land K^+A\) there, and the
same absence on \((s,c)\).

For a test with an undefined argument require undefined
output, with no point condition. For a defined argument
and defined output, require equality of those values and
add its \(A\) label. For a defined argument and undefined
output, add \(\neg A\) at the argument's class.
These cases exhaust the semantics of a test. Every point
condition introduced so far lies at \(x\) or at an
internal class, never at \(y\). In particular, if an
infimum lower bound were \(y\), no possible upper bound
would exceed it and only the invalid-bound rule would apply.

Introduce one fresh existential variable for each equality
class strictly inside \((x,y)\), in its prescribed order.
Distinct classes get distinct variables and a single class
gets only one. Conjoin all point requirements at each class.
At the class of \(x\), their conjunction is the left
label \(\alpha\); an empty list gives true. At each
internal class their conjunction is its point label.

It remains to compile each absence requirement, which has
the form
\[
                    \forall u\,(s<u<t\to L(u))
\]
with \(L=\neg A\) and endpoints already among the
ordered classes. Put \(L\) on every merged open gap
between those endpoint classes, and on every internal
point class strictly between them. Do not put it at
either endpoint. Each point of \((s,t)\) is either
one of those finitely many internal classes or belongs
to exactly one of the indicated open gaps. Conversely
every such class and gap lies in \((s,t)\). The
original universal condition is therefore equivalent
to the conjunction of these transferred conditions.
If there are no classes between adjacent bounds,
the single open gap still receives \(L\); it may
be empty in a particular model. Conjoin all transferred
requirements with the other requirements at each
position. An empty resulting gap list gives true.

The result for this case is \(\alpha(x)\land D(x,y)\).
All its labels are finite Boolean combinations of the
supplied future labels and \(K^+\) applied to such labels,
so they have the actual \(L(U)\) expressions already given.
Only their free current-point argument is copied to a
position. Bound variables inside those copies are made
fresh. The order variables occur in guards and positions,
not as extra parameters inside a label.

We verify the exact meaning of the case formula. An actual
evaluation of the event list in a complete model has
unique definedness statuses and a unique weak order of
its defined values. The infimum-descriptor lemma and
the test semantics verify all constraints in its case.
Assigning the internal class variables those actual
values makes the compiled formula true.

Conversely, suppose a compiled case formula has witnesses
for its internal classes. Restore every point and
absence condition from the gap-transfer equivalence.
Induct through the original event list. The roots
have their assigned values. Invalid or undefined
arguments force the specified undefined result.
For valid infimum bounds, the descriptor identifies
the uniquely correct infimum and status, or identifies
the occurrence set as empty. For a test, its restored
point label gives the specified truth or falsehood
at the already identified argument. Hence every
event's actual value and status agree with the case.
The values also have exactly its ordered partition,
by the strict order of the class witnesses. The case
formula therefore holds exactly when that case is
realized, even if a proposed combination of semantic
labels is inconsistent and makes the formula unsatisfiable.

Taking the finite disjunction of the requested cases
proves the theorem. If no case is requested the output
is false. If all are requested their union is true
under \(x<y\), because the semantics always supplies
one case. Ties, events returning the initial endpoint,
missing events, empty gaps and every Boolean negation
of a status condition are included in the enumeration.
The proof compiles a given list; it has not asserted
that such a list decides every input decomposition.
\end{proof}
\end{annexitem}

\begin{annexitem}{annex:p9.gpss.complete.positive.chain}{thm:p09-gpss-complete-positive-chain}{Infimum selection with a bounded suffix and all nonattainment cases}
\begin{proof}
We first establish the finite splice used in both directions.
Let \(W=(\Theta_1,\Phi_1,\Theta_2,\ldots,\Phi_{m-1},\Theta_m)\).
For \(m\geq2\), suppose a point \(u_1\) satisfies \(E^+_W\), and a
suffix realization
\[
                         u_1<v_2<\cdots<v_m<b
\]
has the remaining point labels and the gaps
\(\Phi_2,\ldots,\Phi_{m-1}\). Choose a provisional
whole-word realization
\(u_1<u_2<\cdots<u_m\), using \(E^+_W\).
If \(u_i<v_i\) for all \(2\leq i\leq m\), it
already ends below \(b\). Otherwise take the least
\(k\in\{2,\ldots,m\}\) with \(v_k\leq u_k\).
For \(k=2\), the assumed \(u_1<v_2\) supplies
the ordering of the new gap. For \(k>2\),
minimality supplies \(u_{k-1}<v_{k-1}<v_k\).
Thus
\[
              u_1,\ldots,u_{k-1},v_k,\ldots,v_m
\]
is a strictly increasing whole-word realization
ending before \(b\): its new gap is contained
in \((u_{k-1},u_k)\) and so has \(\Phi_{k-1}\).
Every other point and gap is inherited at its
own position. Equality \(v_k=u_k\), the first
crossing \(k=2\), and the last crossing \(k=m\)
are all included. For \(m=1\), there is no
suffix and the one point \(u_1<b\) is already
the required realization. This proves the
splice in all applicable lengths.

We now prove correctness of the event recursion
by induction on \(m\). At \(m=0\) both the
target and the Boolean expression impose no
condition. For \(m\geq1\), fix actual \(a<b\)
and assume the assertion for shorter point
lists. The infimum \(c\), if defined, is
the infimum of
\[
                   S=\{u:a<u<b,\ E^+_W(u)\}.
\]
The event \(h\) is undefined precisely when
there is no \(\neg\Phi_0\)-point in \((a,c)\),
when \(c>a\). If \(c=a\), its bounds are
equal and it is undefined by definition;
the corresponding open-gap requirement
is empty. We use these two facts explicitly
in the cases below.

Suppose first that a target chain
\(a<z_1<\cdots<z_m<b\) exists. Its first
point satisfies \(E^+_W\), so \(S\ne\varnothing\).
Completeness supplies \(c=\inf S\) and
\(a\leq c\leq z_1<b\). The initial
\(\Phi_0\)-gap of the chain covers
\((a,c)\), so \(h\) is undefined.
If \(m>1\), its suffix
\(z_2<\cdots<z_m\) lies strictly after
\(c\), has all the shortened list's
point labels and internal gaps, and
satisfies its true initial gap.
The induction hypothesis gives \(R\).
When \(m=1\), \(R\) is true by definition.

If \(c\) has attained status, these
facts suffice for the recurrence.
If it has initial or unattained status,
then \(c<z_1\): the point \(z_1\)
is a member of \(S\), and an unattained
infimum is not that member. Hence
\((c,z_1)\) is a \(\Phi_0\)-neighborhood,
giving \(N^+\Phi_0(c)\). In the
interior unattained case \(a<c<z_1\)
also places \(c\) in the original
initial gap, giving \(\Phi_0(c)\).
In the initial case \(c=a\), that
point is excluded from the original
gap and no \(\Phi_0(c)\) test is
required. Thus the appropriate
test \(j\) or \(k\) is defined,
and every factor of the recurrence
holds.

Conversely suppose the recurrence
holds. Then \(c\) is defined,
the initial gap \((a,c)\) has
\(\Phi_0\) when nonempty, and,
if \(m>1\), the induction
hypothesis supplies a suffix
realization
\[
                         c<v_2<\cdots<v_m<b
\]
with its required internal gaps.
There are three statuses to consider.

In the attained case \(a<c<b\)
and \(E^+_W(c)\) holds. Set
\(u_1=c\). When \(m>1\), apply
the finite splice to this point
and the chosen suffix, obtaining
a whole-word realization inside
\((a,b)\) starting at \(c\).
When \(m=1\), \(c\) itself is
the sole point. In both cases
the initial gap is precisely
\((a,c)\), already supplied
by the undefined \(h\).
No gap label is required at
the first point \(c\).

In either nonattained case,
the infimum-descriptor theorem
gives \(K^+(E^+_W)(c)\).
The defined test \(j\) or \(k\)
gives a point \(d>c\) with
\(\Phi_0\) throughout \((c,d)\).
For \(m>1\), put \(e=\min(d,v_2)\);
for \(m=1\), put \(e=\min(d,b)\).
In both cases \(e>c\).
Apply \(K^+(E^+_W)(c)\) at this
cut to choose
\[
                          c<u_1<e,\qquad E^+_W(u_1).
\]
For \(m>1\), \(u_1<v_2\)
permits the finite splice with
the chosen suffix, yielding a
whole-word realization ending
below \(b\). For \(m=1\),
the chosen point itself lies
below \(b\) and has \(\Theta_1\).

It remains to check its initial
gap, since the splice leaves
its first point \(u_1\)
unchanged. If \(c=a\), the
gap \((a,u_1)\) is a subset
of the supplied \((c,d)\).
If \(a<c\), partition it
into \((a,c)\), the single
point \(c\), and \((c,u_1)\).
The undefined \(h\) supplies
\(\Phi_0\) on the first
piece; test \(k\) supplies
\(\Phi_0(c)\); the neighborhood
supplies the last piece.
Thus the whole initial gap
has \(\Phi_0\). This verifies
every target condition and
proves the reverse direction.

The two nonattained cases show
why an infimum cannot simply
be used as though it were
a first occurrence. Their
actual point \(u_1\) is chosen
strictly above it, the finite
suffix controls the upper
bound, and the point test
at an interior \(c\) covers
the part of the stretched
initial gap that neither
open side interval contains.

Each nonempty recursion step
adds the four nodes \(c,h,j,k\)
and makes one call with one
fewer required point. Its
two endpoint arguments are
the old root or a previously
defined event, so the total
list is acyclic with at
most \(4m\) nonroot nodes.
An unchanged value of the
lower endpoint cannot prevent
termination, since the
point-list length strictly
decreases. All event labels
are finite future formulas:
\(E^+_W\), \(\neg\Phi_0\),
\(N^+\Phi_0\), and
\(\Phi_0\land N^+\Phi_0\).
The two tests have an
earlier event as argument
and never the upper root.

The infimum-case compiler
therefore applies to the
complete Boolean expression,
including the condition that
a nonempty call's arguments
are defined and ordered.
It gives finite left-labelled
decompositions for that
expression and, by selecting
the opposite finite cases,
for its relative complement.
The proof treats a chain
with no final-gap condition.
It has not assigned a
future-only label to the
upper endpoint or removed
a final-gap requirement
from a full decomposition.
\end{proof}
\end{annexitem}

\begin{annexitem}{annex:p9.gpss.complete.bounded.obstruction}{thm:p09-gpss-complete-bounded-obstruction}{A bounded complete-order obstruction for the original grammar}
\begin{proof}
We first verify completeness of the displayed order
\(T=\{0,2\}\cup\{1/n:n\geq1\}\) directly.
Every subset is bounded between 0 and 2. For a nonempty
subset \(A\), its supremum is 2 if \(2\in A\);
otherwise, if \(A\) contains any reciprocal point,
the least denominator among those present supplies
its largest point and hence supremum. In the remaining
case \(A=\{0\}\), its supremum is 0.

Its infimum is 0 if \(0\in A\) or if its reciprocal
denominators are unbounded. In the latter case,
every positive point of \(T\) lies above some
reciprocal point of \(A\), proving greatestness
of the lower bound 0. Otherwise its reciprocal
denominators form a bounded subset of the positive
integers and are finite. The positive members of
\(A\), including a possible 2, form a finite
nonempty set and have a smallest member, its
infimum. These cases exhaust all nonempty subsets.
Thus \(T\) is complete. The point 0 has no immediate
successor, since below any \(1/n\) there is
\(1/(n+1)>0\).

Let \(w=(w_1,w_2,\ldots)\in\{0,1\}^{\mathbb N_{\geq1}}\)
encode the variable interpretation by
\(q_0(1/n)=w_n\). Fix \(q_0(0)=q_0(2)=1\)
and every other predicate false throughout \(T\).
Write \(\mathcal M_w\) for this model. A finite
binary prefix \(s\) defines its cylinder
\[
 [s]=\{w:w\text{ begins with }s\}.
\]
Call a set open if it is a union of such cylinders,
and closed if its complement is open. These
definitions suffice for the argument below.

At any fixed interior point \(1/k\), its entire
future including the present is the finite set
\[
                       \{1/k,1/(k-1),\ldots,1,2\},
\]
with the evident interpretation when \(k=1\).
Only the bits \(w_1,\ldots,w_k\) vary on this tail.
The future-locality lemma therefore implies that
the truth of any permitted \(L(U)\) label at
\(1/k\) depends on only those finitely many bits.
Alternatively the same fact follows directly by
induction: all witnesses and intervening points
of a future modality remain in this finite tail.
Its truth set is consequently a union of some
of the finitely many length-\(k\) cylinders.
It is both open and closed.

Consider any one original decomposition at
the fixed endpoints \(0,2\). Fix a possible
strict interior witness tuple. Every point
label at a witness is a closed condition on
\(w\). For every gap and every actual point
\(u\) in that gap, its future gap label is
likewise a closed condition, since \(u\)
is some reciprocal point and has a finite
future. A universal gap condition is the
intersection of all these closed conditions.
Arbitrary intersections of closed sets are
closed: their complementary union is open.
Conjoining the finitely many point and gap
requirements for this fixed tuple is thus
closed. If the tuple has no interior points,
the same reasoning applies to its single
universal gap.

There are only countably many finite tuples
of points of \(T\). Existentially choosing
the witnesses therefore makes the truth
set of one decomposition a countable
union of closed sets. A finite disjunction
of decompositions is still a countable
union of closed sets. No label is evaluated
at 0 in this argument: that exclusion
is exactly the original grammar's
endpoint convention. There are no
cofinal conjuncts for these finite
endpoints.

The particular word \(D_A(0,2)\) has a
witness \(z=1/k\) exactly when
\[
                         w_n=1\quad\text{for all }n>k.
\]
Its final gap has label true and imposes
no other condition. Thus its truth set
is \(\mathrm{FIN}\), the set of binary
sequences with only finitely many zero
entries. Its complement is
\(\mathrm{INF}\), those with infinitely
many zero entries.

We give the full reason that \(\mathrm{INF}\)
is not a countable union of closed sets.
Suppose \(\mathrm{INF}=\bigcup_{j\geq1}C_j\),
where each \(C_j\) is closed. Begin with
the empty prefix. At step \(j\), extend
the current finite prefix by an infinite
tail of ones, giving a sequence \(v\)
with finitely many zeros. It belongs
to none of the \(C_j\), since every
\(C_j\) is contained in \(\mathrm{INF}\).
In particular \(v\notin C_j\).
The complement of the closed set \(C_j\)
is open, so some sufficiently long
initial prefix of \(v\), extending
the current prefix, has its entire
cylinder disjoint from \(C_j\).
Append one zero to that prefix to
form the prefix for the next step.
Its cylinder is still disjoint from
\(C_j\).

The nested finite prefixes determine
an infinite sequence \(w\), since
their lengths strictly increase.
A new zero was appended at every
step, so \(w\in\mathrm{INF}\).
For every \(j\), all later prefixes
stay in the cylinder chosen to avoid
\(C_j\). Hence \(w\notin C_j\)
for every \(j\), a contradiction.
This proves the assertion without
using a descriptive-set-theoretic
classification theorem as an
unproved premise.

If a finite original decomposition
disjunction expressed the guarded
complement of \(D_A\) on complete
orders, it would do so at these
fixed actual endpoints for every
\(\mathcal M_w\). Its truth set
would be a countable union of
closed sets, equal to
\(\mathrm{INF}\), which has just
been proved impossible. This
establishes the claimed obstruction.
It does not contradict a
future-only expression at 0,
whose future includes all the
varying reciprocal points.
\end{proof}
\end{annexitem}

\begin{annexitem}{annex:p9.gpss.initial.gap.complement}{thm:p09-gpss-initial-gap-complement}{The exact endpoint-labelled correction for an initial homogeneous gap}
\begin{proof}
Fix actual \(x<y\). If these endpoints are adjacent,
there is no possible interior witness and \(D_A(x,y)\)
is false. At such an \(x\), the immediate successor
\(y\) makes \(K^+(\neg A)(x)\) false, regardless
of all labels. The first displayed alternative
therefore gives exactly this case.

Suppose instead that \((x,y)\) contains a point.
We prove
\[
                       D_A(x,y)\longleftrightarrow N^+A(x).
 \tag{1}
\]
A witness for \(D_A\) immediately witnesses
the neighborhood formula. Conversely let \(d>x\)
witness \(N^+A(x)\). If \(d<y\), choose \(z=d\).
Its initial gap satisfies \(A\), so it witnesses
\(D_A\). If \(d\geq y\), choose any point
\(z\in(x,y)\), which exists by the present case
assumption. Then \(z<d\) and
\((x,z)\subseteq(x,d)\), again supplying the
required initial gap. This also covers equality
\(d=y\). The final gap has label true in both cases.
Thus (1) holds.

Negating \(N^+A\) gives
\(\partial^+(\neg A)\) by its explicit
strict-until definition. Since \(y>x\)
already supplies \(F\top\), this is
equivalent to \(K^+(\neg A)(x)\).
Together with the adjacent case it proves
the body's equivalence under \(x<y\).

The zero-point decomposition with false
gap says that \((x,y)\) is empty and
therefore represents the adjacent case.
The zero-point true-gap decomposition
is true on every ordered endpoint pair.
Conjoining its left label
\(K^+(\neg A)\) represents the other
case. Neither construction tests
\(A\) at an external endpoint.

On the reciprocal-point countermodel,
\(K^+(\neg q_0)(0)\) says that
every right neighborhood of 0
contains a zero-labelled reciprocal
point. This is exactly infinitely
many zero entries. The corrected
left-labelled formula therefore
expresses the complement whose
representation in the original
class the preceding theorem rules
out.
\end{proof}
\end{annexitem}
